\begin{table}[htbp]
\centering
\caption{虚拟台风情景灵敏度解释}
\label{tab:problem3_sensitivity}
\begin{tabular}{c l l r p{7.2cm}}
\hline
情景 & 主控因素 & 最敏感指标 & 观测变化 & 物理解释 \\
\hline
S1 & 强度增强 & $A_{10}$ & 6.70\% & 风速增强并伴随气压降低，主要抬升强降水面积和雨强尾部，对最大雨强、P99 和强降水覆盖更敏感。 \\
S2 & 路径近岸/西移 & 雨区质心偏移 & 104.76 km & 路径向近岸侧偏移后，storm-relative 雨带投影到固定经纬度网格的位置改变，主要表现为降水落区和重叠率变化。 \\
S3 & 近岸段偏移 & $1-\mathrm{IoU}_{D10}$ & 22.27\% & 局地近岸段扰动改变沿海影响带位置，但不显著改变整个台风生命史内的强度和总雨量条件。 \\
S4 & 移动速度减慢 & 最大累计降水 & 66.98\% & 移动速度降低延长固定地理位置受强降水影响的时间，因此累计降水和固定格点持续时间显著增大。 \\
S5 & 复合高风险 & 最大累计降水 & 30.99\% & 中等强度增强、近岸偏移和轻度减速共同作用，使雨强、落区和持续时间多个风险维度同步增加。 \\
\hline
\end{tabular}
\end{table}
